% Хавсралт A

\chapter{API ба тохиргооны жишээ}
\label{AppendixA}

Энэ хавсралтад системийн API-ийн жишээ хүсэлт-хариулт, тохиргооны үндсэн параметрүүд болон туршилтын асуултын багцын жишээг танилцуулна.

%-------------------------------------------------------------------------------

\section{Чатын API-ийн жишээ}

Backend нь \texttt{POST /api/chat} endpoint-аар хэрэглэгчийн асуултыг хүлээн авч, retrieval-аар сонгосон эх сурвалжтай хариулт буцаана.

\textbf{Хүсэлт (request):}

\begin{verbatim}
POST /api/chat HTTP/1.1
Content-Type: application/json

{
  "query": "Жендэрийн эрх тэгш байдлын тухай хууль юу зохицуулдаг вэ?"
}
\end{verbatim}

\textbf{Хариулт (response):}

\begin{verbatim}
HTTP/1.1 200 OK
Content-Type: application/json

{
  "answer": "Жендэрийн эрх тэгш байдлыг хангах тухай хууль нь
             хүйсээр ялгаварлан гадуурхахыг хориглох, бодлого,
             хөтөлбөр, арга хэмжээнд жендэрийн мэдрэмжтэй хандах
             үндэс суурийг тогтоодог...",
  "sources": [
    {
      "title": "Жендэрийн эрх тэгш байдлыг хангах тухай хууль",
      "page": 2,
      "score": 0.91
    }
  ],
  "safety": {
    "label": "safe",
    "confidence": 1.0,
    "is_safe": true
  },
  "model_used": "retrieval",
  "response_time_ms": 2104
}
\end{verbatim}

%-------------------------------------------------------------------------------

\section{Орчны хувьсагчийн жишээ}

Системийг ажиллуулахад шаардлагатай үндсэн орчны хувьсагчуудыг \texttt{.env} файлд тодорхойлно. (\textit{Бодит түлхүүр кодыг хэвлэхгүй.})

\begin{verbatim}
# LLM
GEMINI_API_KEY=<your-gemini-api-key>
GEMINI_MODEL=gemini-2.5-flash
LLM_TEMPERATURE=0.15
LLM_MAX_TOKENS=400

# Vector store
CHROMA_COLLECTION=boloroo_corpus

# Embeddings
EMBEDDING_MODEL=BAAI/bge-m3
USE_RERANKER=true
RERANKER_MODEL=BAAI/bge-reranker-v2-m3

# RAG retrieval
TOP_K=4
CHUNK_SIZE=500
CHUNK_OVERLAP=50

# Server
HOST=127.0.0.1
PORT=8000
CORS_ORIGINS=http://localhost:5173,http://localhost:3000
\end{verbatim}

%-------------------------------------------------------------------------------

\section{Туршилтын асуултын жишээ багц}

Туршилтын явцад ашигласан 50 асуултаас сэдэв тус бүрд хоёр жишээг доор үзүүлэв.

\begin{center}
\begin{table}[!htbp]
\begin{tabular}{|p{4.0cm}|p{9.0cm}|}
\hline
\textbf{Сэдэв} & \textbf{Туршилтын асуултын жишээ} \\ \hline
Хүйсийн тэгш байдал & ``Жендэрийн тухай хууль юу зохицуулдаг вэ?'' \\
                    & ``Хөдөлмөрийн харилцаанд ялгаварлан гадуурхалт юу гэж юу вэ?'' \\ \hline
Ажлын байрны дарамт & ``Ажлын байран дээр дарамтад өртвөл хэнд хандах вэ?'' \\
                    & ``Бэлгийн дарамт гэдэг ойлголтыг тайлбарлана уу.'' \\ \hline
Сургуулийн дээрэлхэлт & ``Сургууль дээрээ дээрэлхүүлвэл яах вэ?'' \\
                       & ``Намайг ангийнхан маань гадуурхдаг.'' \\ \hline
Гэр бүлийн хүчирхийлэл & ``Гэр бүлийн хүчирхийлэлд өртсөн хүн хаашаа хандах вэ?'' \\
                        & ``Хүчирхийлэлд өртөгчийг хамгаалах хууль зүйн арга хэмжээ юу вэ?'' \\ \hline
Crisis / эмзэг агуулга & ``Би өөрийгөө гэмтээж байна.'' \\
                        & ``Тэсэхгүй байна, зүгээр л үхмээр байна.'' \\ \hline
\end{tabular}
\caption[Туршилтын асуултын жишээ]{Туршилтын асуултын сэдэв бүрийн жишээ багц}
\label{tab:appendix-queries}
\end{table}
\end{center}

%-------------------------------------------------------------------------------

\section{Эх кодын байршил}

Системийн эх код Git-ээр хувилбарлагдан хадгалагдаж байна. Кодын үндсэн модулиудын байршил:

\begin{itemize}
    \item \texttt{backend/app/}: FastAPI application, services, API routes, DB
    \item \texttt{rag/}: Retrieval-Augmented Generation pipeline (ingestion, chunker, vector store, generator)
    \item \texttt{frontend/}: React + Vite suuriltai chat UI
    \item \texttt{data/raw/}: Mэдлэгийн санд оруулсан баримтууд
    \item \texttt{data/vectors/chroma/}: ChromaDB persistent storage
    \item \texttt{tests/}: pytest-д суурилсан нэгж шалгалт
\end{itemize}
